\chapter{Segitiga kongruen}\label{chap:cong}

\section{Sisi-sudut-sisi}

Tujuan kita selanjutnya adalah memberikan syarat-syarat yang menjamin kekongruenan dua segitiga.

Salah satu syarat tersebut diberikan dalam Aksioma~\ref{def:birkhoff-axioms:3}; aksioma itu menyatakan bahwa jika dua pasang sisi dari dua segitiga sama panjang, dan sudut-sudut apitnya sama hingga tanda, maka kedua segitiga tersebut kongruen.
Aksioma ini juga disebut {}\emph{syarat kekongruenan sisi-sudut-sisi}, atau singkatnya, \index{syarat kekongruenan SAS}\emph{SAS}.

\section{Sudut-sisi-sudut}

\begin{thm}[\abs]{Syarat ASA}\label{thm:ASA}\index{syarat kekongruenan ASA}
Andaikan bahwa
\begin{align*}
AB&=A'B',
&
\measuredangle A B C &= \pm\measuredangle A' B' C',
&
\measuredangle C A B&=\pm\measuredangle C' A' B'
\end{align*}
 dan $\triangle A' B' C'$ tak degenerat.
Maka
$$\triangle A B C\cong\triangle A' B' C'.$$

\end{thm}

Untuk segitiga degenerat, pernyataan ini tidak berlaku.
Sebagai contoh, tinjau sebuah segitiga dengan panjang sisi $1$, $4$, $5$
dan segitiga lain dengan panjang sisi $2$, $3$,~$5$.

\parit{Bukti.}
Menurut Teorema~\ref{thm:signs-of-triug},
berlaku salah satu dari dua kemungkinan berikut:
\[
\measuredangle A B C = \measuredangle A' B' C'
\quad\text{dan}\quad
\measuredangle C A B=\measuredangle C' A' B'
\eqlbl{eq:+angles}\]
atau
\[
\measuredangle A B C = -\measuredangle A' B' C'
\quad\text{dan}\quad
\measuredangle C A B=-\measuredangle C' A' B'.
\eqlbl{eq:-angles}\]
Andaikan \ref{eq:+angles} berlaku;
kasus \ref{eq:-angles} serupa.

\begin{wrapfigure}{o}{26mm}
\vskip-2mm
\centering
\includegraphics{mppics/pic-36}
\end{wrapfigure}

Misalkan $C''$ adalah titik pada sinar $[A' C')$ sedemikian sehingga $A' C''\z=A C$.

Menurut Aksioma~\ref{def:birkhoff-axioms:3},
$\triangle A' B' C''\cong \triangle A B C$.
Dengan menerapkan kembali Aksioma~\ref{def:birkhoff-axioms:3},
kita memperoleh
$$\measuredangle A' B' C'' = \measuredangle A B C=\measuredangle A' B' C'.$$
Menurut Aksioma~\ref{def:birkhoff-axioms:2a}, $[B' C')=[B' C'')$.
Dengan demikian,
$C''$ terletak pada $(B' C')$ maupun~$(A' C')$.

Karena $\triangle A' B' C'$ tak degenerat, $(A' C')$ berbeda dari~$(B' C')$.
Dengan menerapkan Aksioma~\ref{def:birkhoff-axioms:1}, kita memperoleh $C''=C'$.

Oleh karena itu,
$\triangle A' B' C'=\triangle A' B' C''\cong\triangle A B C$.
\qeds

\section{Segitiga sama kaki}

Segitiga yang memiliki dua sisi sama panjang disebut \index{segitiga sama kaki}\emph{sama kaki};
sisi yang tersisa disebut \index{alas}\emph{alas}.


\begin{thm}[\abs]{Teorema}\label{thm:isos}
Andaikan $\triangle A B C$ adalah segitiga sama kaki dengan alas $[A B]$.
Maka
$$\measuredangle A B C\equiv -\measuredangle B A C.$$
Selain itu, kebalikannya berlaku jika $\triangle A B C$ tak degenerat.
\end{thm}

Bukti berikut berasal dari Pappus dari Aleksandria.



\parit{Bukti.}
Perhatikan bahwa
$$
\measuredangle A C B \equiv -\measuredangle B C A,
\quad
C A = C B,
\quad \text{dan oleh karena itu}\quad
C B=C A.$$

\begin{wrapfigure}[9]{r}{25mm}
\vskip-4mm
\centering
\includegraphics{mppics/pic-38}
\end{wrapfigure}

Menurut Aksioma~\ref{def:birkhoff-axioms:3},
$$\triangle C A B\cong\triangle C B A.$$
Dengan menerapkan teorema tentang tanda sudut-sudut segitiga (\ref{thm:signs-of-triug}) dan Aksioma~\ref{def:birkhoff-axioms:3} sekali lagi,
kita memperoleh
$$\measuredangle B A C
\equiv -\measuredangle A B C.$$

Untuk membuktikan kebalikannya, andaikan bahwa
$\measuredangle C A B \z\equiv - \measuredangle C B A$.
Menurut syarat ASA (\ref{thm:ASA}), $\triangle C A B\z\cong\triangle CBA$.
Oleh karena itu,~$C A\z=C B$.
\qeds

Segitiga yang memiliki tiga sisi sama panjang disebut \index{segitiga sama sisi}\emph{sama sisi}.

\begin{thm}{Latihan}\label{ex:equilateral}
Misalkan $\triangle ABC$ adalah segitiga sama sisi.
Tunjukkan bahwa
\[\measuredangle ABC=\measuredangle BCA=\measuredangle CAB.\]

\end{thm}


\section{Sisi-sisi-sisi}

\begin{thm}[\abs]{Syarat SSS}\label{thm:SSS}\index{syarat kekongruenan SSS}
$\triangle A B C\cong\triangle A' B' C'$ jika
$$A' B'=A B,
\quad
B' C'=B C,
\quad
\text{dan}
\quad
C' A'=C A.$$

\end{thm}

Perhatikan bahwa syarat ini juga berlaku untuk segitiga degenerat.

\parit{Bukti.}
Pilih $C''$ sedemikian sehingga $A' C''= A' C'$ dan $\measuredangle B' A' C''= \measuredangle B A C$.
Menurut Aksioma~\ref{def:birkhoff-axioms:3},
$$\triangle A' B' C''\cong\triangle A B C.$$

Cukup dibuktikan bahwa
$$\triangle A' B' C'\cong\triangle A' B' C''.\eqlbl{eq:A'B'C'simA'B'C''}$$
Syarat \ref{eq:A'B'C'simA'B'C''} jelas berlaku jika $C''\z=C'$.
Dengan demikian, tinggal meninjau kasus $C''\z\ne C'$.

\begin{wrapfigure}{o}{30mm}
\centering
\includegraphics{mppics/pic-40}
\end{wrapfigure}

Jelas bahwa sisi-sisi yang bersesuaian dari $\triangle A' B' C'$ dan $\triangle A' B' C''$ sama panjang.
Oleh karena itu, $\triangle C' A' C''$ dan $\triangle C' B' C''$ merupakan segitiga sama kaki.
Menurut Teorema~\ref{thm:isos}, berlaku
\begin{align*}
 \measuredangle A' C'' C'&\equiv -\measuredangle A' C' C'',
\\
\measuredangle C' C'' B'&\equiv -\measuredangle C'' C' B'.
\end{align*}
Dengan menjumlahkan keduanya, kita memperoleh
$$\measuredangle A' C'' B'
\equiv -\measuredangle A' C' B'.$$
Dengan menerapkan kembali Aksioma~\ref{def:birkhoff-axioms:3},
kita memperoleh \ref{eq:A'B'C'simA'B'C''}.
\qeds

\begin{thm}[\abs]{Akibat}\label{cor:degenerate-trig}
Jika $AB+BC=AC$, maka $B\in [AC]$.
\end{thm}

\parit{Bukti.}
Karena $AB+BC=AC$, kita dapat memilih $B'\in [AC]$ sedemikian sehingga $AB=AB'$;
perhatikan bahwa $BC\z=B'C$.

Kita dapat mengandaikan bahwa $AB>0$ dan $BC>0$;
jika tidak, $A=B$ atau $B=C$, dan pernyataan langsung berlaku.
Dalam kasus ini, $B'$ terletak di antara $A$ dan $C$.
Menurut Teorema~\ref{thm:straight-angle}, $\measuredangle AB'C=\pi$.

Menurut SSS,
\[\triangle ABC\cong \triangle AB'C.\]
Oleh karena itu, $\measuredangle ABC=\pi$.
Menurut Teorema~\ref{thm:straight-angle}, $B$ terletak di antara $A$ dan $C$.
\qeds



\begin{thm}{Latihan lanjutan}\label{ex:SMS}
Tinjau dua segitiga $A B C$ dan $A' B' C'$.
Misalkan $M$ adalah titik tengah $[A B]$, dan
$M'$ adalah titik tengah $[A' B']$.
Andaikan $C' A'=C A$, $C' B'= C B$, dan $C' M'\z= C M$.
Buktikan bahwa
\[\triangle A' B' C'\cong\triangle A B C.\]

\end{thm}

{

\begin{wrapfigure}[6]{r}{26mm}
\vskip-0mm
\centering
\includegraphics{mppics/pic-42}
\end{wrapfigure}

\begin{thm}{Latihan}\label{ex:isos-sides}
Misalkan $\triangle A B C$ adalah segitiga sama kaki dengan alas $[A B]$.
Andaikan titik $A'$ terletak di antara $B$ dan $C$,
titik $B'$ terletak di antara $A$ dan $C$,
dan $C A'\z=C B'$.
Tunjukkan bahwa
\begin{enumerate}[(a)]
\item $\triangle A A' C\cong \triangle B B' C$;
\item $\triangle A B B'\cong \triangle B A A'$.
\end{enumerate}

\end{thm}

\begin{thm}[!]{Latihan}\label{ex:ABC-motion}
Misalkan $\triangle ABC$ adalah segitiga tak degenerat, dan
$f$ adalah pemetaan yang mempertahankan jarak dari bidang ke dirinya sendiri
sedemikian sehingga
$$f(A)=A,
\quad
f(B)=B,
\quad
\text{dan}
\quad
f(C)=C.$$

Tunjukkan bahwa $f$ adalah pemetaan identitas;
artinya, $f(X)=X$ untuk setiap titik $X$ pada bidang.
\end{thm}

Pernyataan dalam latihan tersebut disebut \emph{lema tiga paku}, karena tiga paku menahan bidang pada tempatnya. %(Tidak seperti di Yerusalem, tidak ada yang terluka di sini.)

}

%{

%\begin{wrapfigure}{r}{28mm}\vskip-11mm\centering\includegraphics{mppics/pic-43}\end{wrapfigure}

%\begin{thm}{Latihan}\label{ex:3-isos}
%Andaikan bahwa $\triangle ABC'$, $\triangle BCA'$, dan $\triangle CAB'$ adalah segitiga sama kaki yang terletak di luar $\triangle ABC$;
%artinya,
%$A$ dan $A'$ terletak pada sisi-sisi yang berlawanan dari $(BC)$,
%$B$ dan $B'$ terletak pada sisi-sisi yang berlawanan dari $(CA)$,
%$C$ dan $C'$ terletak pada sisi-sisi yang berlawanan dari $(AB)$.

%Tunjukkan bahwa $AA'=BB'=CC'$.
%\end{thm}

%}

\section{Tentang sisi-sisi-sudut dan sisi-sudut-sudut}

Dalam setiap syarat SAS, ASA, dan SSS, kita menetapkan tiga bagian yang bersesuaian dari kedua segitiga.
Sekarang, mari kita bahas tiga bagian bersesuaian lainnya.

Kombinasi pertama disebut {}\emph{sisi-sisi-sudut}, atau singkatnya SSA;
kombinasi ini menetapkan dua sisi dan satu sudut yang bukan sudut apit.
Syarat ini tidak cukup untuk menjamin kekongruenan.
Dengan kata lain, terdapat dua segitiga tak degenerat $ABC$ dan $A'B'C'$ sedemikian sehingga
\[AB=A'B',\quad BC=B'C',\quad \measuredangle BAC\equiv\pm \measuredangle B'A'C',\]
tetapi $\triangle ABC\not\cong\triangle A'B'C'$ dan $AC\ne A'C'$.

\begin{wrapfigure}{r}{35mm}
\vskip-6mm
\centering
\includegraphics{mppics/pic-44}
\end{wrapfigure}

Kita tidak akan menggunakan pernyataan negatif ini, sehingga tidak perlu membuktikannya secara formal.
Anda dapat menduga sebuah contoh dari gambar.

Kombinasi kedua disebut {}\emph{sisi-sudut-sudut}, atau singkatnya SAA;
kombinasi ini menetapkan satu sisi dan dua sudut, salah satunya berhadapan dengan sisi tersebut.
Ini memberikan suatu syarat kekongruenan;
artinya, jika salah satu dari segitiga $ABC$ atau $A'B'C'$ tak degenerat dan
$AB=A'B'$, $\measuredangle ABC\equiv\pm \measuredangle A'B'C'$, $\measuredangle BCA\z\equiv\pm \measuredangle B'C'A'$,
maka $\triangle ABC\cong\triangle A'B'C'$.

Syarat SAA tidak akan digunakan secara langsung selanjutnya.
Salah satu bukti untuk syarat ini dapat diperoleh dari ASA dan teorema tentang jumlah sudut segitiga, yang dibuktikan kemudian (lihat~\ref{thm:3sum}).
Untuk bukti yang lebih langsung, lihat Latihan~\ref{ex:SAA}.

Kombinasi lainnya disebut {}\emph{sudut-sudut-sudut}, atau singkatnya AAA.
Menurut Aksioma~\ref{def:birkhoff-axioms:4}, syarat ini bukan syarat kekongruenan pada bidang Euklides.
Namun, pada bidang hiperbolik syarat ini merupakan syarat kekongruenan; lihat \ref{thm:AAA}.

{

\section{Konstruksi}

Masalah konstruksi merupakan sumber latihan geometri yang berharga,
yang akan kita gunakan lebih lanjut dalam buku ini.
Bab~\ref{chap:car} akan dikhususkan untuk topik tersebut.

\index{konstruksi dengan penggaris dan jangka}\emph{Konstruksi dengan penggaris dan jangka} adalah konstruksi yang menggunakan penggaris dan jangka ideal.
Penggaris ideal hanya dapat digunakan untuk menggambar garis yang melalui dua titik tertentu.
Jangka ideal hanya dapat digunakan untuk menggambar lingkaran dengan pusat dan jari-jari tertentu.
Artinya, jika diberikan tiga titik $A$, $B$, dan $O$,
kita dapat menggambar lingkaran berjari-jari $AB$ yang berpusat di~$O$.
Kita juga dapat menandai titik-titik baru pada bidang,
serta pada garis, lingkaran,
dan titik perpotongan yang telah dikonstruksi (dengan asumsi titik-titik tersebut ada).

Kita juga dapat meninjau perangkat konstruksi yang berbeda.
Sebagai contoh,
kita dapat menggunakan penggaris saja,
atau menciptakan alat baru,
misalnya alat yang menghasilkan titik tengah dari dua titik tertentu,
atau cerminan satu titik terhadap titik lainnya.

\begin{thm}{Masalah}
Diberikan dua garis $\ell$ dan $m$ serta sebuah titik $M$ yang tidak terletak pada salah satu pun dari kedua garis tersebut,
konstruksikan ruas garis yang ujung-ujungnya masing-masing terletak pada $\ell$ dan $m$, serta yang dibagi dua oleh~$M$.
\end{thm}

Tahap-tahap konstruksi berikut ditunjukkan dalam gambar.

\begin{center}
\includegraphics[width=.29\textwidth]{mppics/pic-445}\hfill
\includegraphics[width=.29\textwidth]{mppics/pic-446}\hfill
\includegraphics[width=.29\textwidth]{mppics/pic-447}

\medskip

\includegraphics[width=.29\textwidth]{mppics/pic-448}\hspace{.08\textwidth}
\includegraphics[width=.29\textwidth]{mppics/pic-449}
\end{center}

\parit{Konstruksi.}
\begin{enumerate}[1.]
\item Gambarlah sebuah lingkaran berpusat di $M$ yang memotong garis $\ell$ di dua titik, misalnya $P$ dan $Q$.
\item Gambarlah garis $(PM)$ dan tandai titik potongnya yang lain dengan lingkaran sebagai $P'$.
Dengan cara serupa, gambarlah $(QM)$ dan tandai titik potongnya yang lain sebagai $Q'$.
\item Gambarlah garis $\ell'=(P'Q')$ dan tandai sebagai $A$ titik potongnya dengan $m$ (jika ada).
\end{enumerate}

}

\begin{enumerate}[1.]\setcounter{enumi}{3}
\item Gambarlah garis $(MA)$;
tandai sebagai $B$ titik potongnya dengan $\ell$.
Maka $[AB]$ adalah ruas yang diminta.
\end{enumerate}


Sering kali perlu dilakukan \emph{analisis masalah}:
membuktikan bahwa hasil konstruksi memenuhi syarat-syarat yang diminta, serta menentukan kapan konstruksi dapat dilakukan dan berapa banyak solusi yang mungkin ada.
Mari kita lakukan analisis tersebut untuk contoh di atas.


\parit{Analisis masalah.}
Perhatikan bahwa $P'$ dan $Q'$ masing-masing merupakan cerminan $P$ dan $Q$ terhadap~$M$.
Oleh karena itu, garis $\ell'=(P'Q')$ adalah cerminan $\ell=(PQ)$ terhadap~$M$.

Karena $A$ terletak pada $\ell'$,
titik $B$ adalah cerminan $A$ terhadap $M$.
Memang, $B$ adalah titik potong (yang pasti tunggal) antara $(MA)$ dan $\ell$, sedangkan $\ell'$ merupakan cerminan $\ell$.
Menurut definisi pencerminan, $M$ adalah titik tengah $[AB]$.

Garis $\ell'$ mungkin tidak memiliki titik potong dengan $m$.
Dalam kasus ini, masalah tersebut tidak memiliki solusi.
Mungkin juga terjadi bahwa $\ell'=m$;
dalam kasus ini, terdapat tak berhingga banyak solusi --- sembarang titik pada $m$ dapat ditandai sebagai~$A$.
Dalam semua kasus selebihnya, terdapat tepat satu solusi.
\qeds

\begin{thm}{Latihan}\label{ex:motion}
Andaikan $\triangle ABC$ dan $\triangle A'B'C'$ adalah segitiga tak degenerat yang kongruen;
khususnya, terdapat transformasi isometrik $f$ sedemikian sehingga
\[f\colon A\mapsto A',\quad f\colon B\mapsto B',\quad\text{dan}\quad f\colon C\mapsto C'.\]

Diberikan sebuah titik $X$, konstruksikan $X'=f(X)$.
\end{thm}
